\scanpage{3}
\deleted{
Normalement, on voudrait que
\diagram{\begin{tikzcd}
\uline{\unsure{Ch}}(C) \arrow[r,"p"] & C\times C
\end{tikzcd}}
\unsure{ayant|ait} des propriétés de type ``fibration''. Or

1) au sens catégorique, ce n'est ni une fibration, ni une cofibration --- et \unsure{$=$|même} les \deleted{projections} composés
\[
\uline{$s$}\circ p,\ \uline{$\unsure{b}$}\circ p\colon\ \uline{\unsure{Ch}}(C)\to C
\]
ne le sont pas.

2) De $\unsure{\simeq}$, \deleted{\unsure{ni}} $p$ \deleted{\illegible{1}} n'est ni coh. propre, ni coh. lisse.
Peut-être au moins \deleted{\unsure{$p$}} $\uline{s}\circ p$ et $\uline{\unsure{b}}\circ p$ le sont-ils \deleted{\illegible{1}} l'un ou l'autre? Mais ce sont plutôt les props de $p$ lui-même qui est important!

3) Il semblerait que $p$ soit une $W_\infty$-fibration.
(``fibration \unsure{cohomologique}'') C'est une chose à vérifier.

\uline{Question} Les $W_\infty$-fibrations, et la notion correspondante de cofibrations \deleted{\unsure{pour}} \inserted{(pour la L.L.P)}, font-ils une \deleted{des} structure de Quillen sur $(\mathrm{Cat})$?
}
\mspage{2}

